
\documentclass{.local/lib/classes/ndvr-vector-image-standalone}

\title[NDVR thesis-abstract]{%
    Темпоральная модель визуального потока видео%
}

\textwidth=40cm

\begin{document}

    %% My Types

    \newcommand{\umldtp}[1]{\textsf{#1}}
    \newcommand{\umlctp}[1]{\textsl{$\star$ #1}}
    \newcommand{\umlfld}[1]{\textsf{#1}}


    %% Tabular Styling

    \newcolumntype{F}[1]{>{\raggedleft\arraybackslash}m{#1}}
    \newcolumntype{T}[1]{>{\raggedright\arraybackslash}m{#1}}

    \setlength\dashlinedash{0.2pt}
    \setlength\dashlinegap{2.5pt}
    \setlength\arrayrulewidth{0.2pt}

    \renewcommand{\arraystretch}{1.2}
    \setlength{\tabcolsep}{0.1em}


    \colorlet{commonClass}{yellow!10}
    \colorlet{commonNote}{green!10}
    \colorlet{commonDraw}{lightgray}

    \colorlet{emphClass}{yellow!20}
    \colorlet{emphDraw}{black}


    \tikzstyle{tikzuml transition style}=[%
        color=\tikzumlDefaultDrawColor,%
        very thick,% 
        rounded corners,% 
        -angle 45%
    ]%


    %% Main Part

    \sf\singlespacing
    \begin{tikzpicture}[
        very thick,
        node distance=8em,
    ]

        \umlstateinitial[y=0,name=Init]

        \begin{umlstate}[
            name=Physical,
            y=-6,
            fill=commonClass,
        ]{Физическое представление видео}
        \end{umlstate}
        \umltrans
            {Init}
            {Physical}

        \begin{umlstate}[
            name=Numerical,
            y=-12,
            fill=commonClass,
        ]{Численное представление видео}
        \end{umlstate}
        \umltrans
            {Physical}
            {Numerical}

        \begin{umlstate}[
            name=Segments,
            y=-18,
            fill=commonClass,
        ]{Сегменты}
            \begin{umlstate}[
                name=Key-frames,
                x=-5,
            ]{Ключевые кадры}
            \end{umlstate}
            \begin{umlstate}[
                name=Shots,
                x=0,
            ]{Cъёмки}
            \end{umlstate}
            \begin{umlstate}[
                name=Scenes,
                x=5, 
            ]{Сцены}
            \end{umlstate}
        \end{umlstate}
        \umltrans[
            arg1={Сегментация}
        ]
            {Numerical}
            {Segments}

        \begin{umlstate}[
            name=Features,
            fill=commonClass,
            y=-24
        ]{Сегменты}
            \begin{umlstate}[
                name=Frames1,
                x=-5,
            ]{Кадры 1}
            \end{umlstate}
            \begin{umlstate}[
                name=Frames,
                x=0,
            ]{Кадры 1}
            \end{umlstate}
            \begin{umlstate}[
                name=ZFrames,
                x=5, 
            ]{Кадры 2}
            \end{umlstate}
        \end{umlstate}
        \umltrans[
            arg={Выборка характеристик}
        ]
            {Segments}
            {Features}

        \begin{umlstate}[
            name=Information,
            fill=commonClass,
            y=-30
        ]{Сегменты}
            \begin{umlstate}[
                name=Frames1,
                x=-5,
            ]{Кадры 1}
            \end{umlstate}
            \begin{umlstate}[
                name=Frames,
                x=0,
            ]{Кадры 1}
            \end{umlstate}
            \begin{umlstate}[
                name=ZFrames,
                x=5, 
            ]{Кадры 2}
            \end{umlstate}
        \end{umlstate}
        \umltrans[
            arg={Извлечение информации}
        ]
            {Features}
            {Information}
            
        % \umlstateinitial[name=init]
        % \begin{umlstate}[
        %     name=Video, 
        %     x=0, y=0,
        % ]{Видео}
        % \end{umlstate}

        % % \begin{umlstate}[
        % %     name=Features, 
        % %     x=3, y=0,
        % %     ]{Характеристики}
        % % \end{umlstate}
        % % \begin{umlstate}[
        % %     name=Information, 
        % %     x=4, y=0,
        % % ]{Информация}
        % % \end{umlstate}


    % \begin{umlstate}[name=Amain]{Etat global de l'objet A}
    %     \begin{umlstate}[name=Bgraph, fill=red!20]{graphe B}
    %         \umlstateinitial[name=Binit]
    %         \umlbasicstate[y=-4, name=test1, fill=white]{test1}
    %         \umltrans{Binit}{test1}
    %         \umltrans[recursive=20|60|2.5cm, recursive direction=right to top, arg={op1}, pos=1.5]{test1}{test1}
    %         \umltrans[recursive=160|120|2.5cm, recursive direction=left to top, arg={op2}, pos=1.5]{test1}{test1}
    %         \umltrans[recursive=-160|-120|2.5cm, recursive direction=left to bottom, arg={op3}, pos=1.5]{test1}{test1}
    %         \umltrans[recursive=-20|-60|2.5cm, recursive direction=right to bottom, arg={op4}, pos=1.5]{test1}{test1}
    %         \umlbasicstate[y=-8, name=test2, fill=white]{test2}
    %         \umltrans[recursive=-160|-120|2.5cm, recursive direction=left to bottom, arg={op5}, pos=1.5]{test2}{test2}
    %         \umltrans{test1}{test2}
    %         \umlstatefinal[x=3, y=-7.75, name=Bfinal]
    %         \umltrans{test2}{Bfinal}
    %     \end{umlstate}
    %     \umlstateinitial[x=6, y=1, name=Ainit]
    %     \umlVHtrans[anchor2=40]{Ainit}{Bgraph}
    %     \umlstatefinal[x=6, y=-3.5, name=Afinal]
    %     \umlHVtrans[anchor1=30]{Bgraph}{Afinal}
    %     \umlbasicstate[x=6, y=-6, name=visu, fill=green!20]{Visualisation}
    %     \umlHVtrans{Bfinal}{visu}
    %     \umltrans{visu}{Afinal}
    %     \umltrans[recursive=-20|-60|2.5cm, recursive direction=right to bottom, arg=a, pos=1.5]{visu}{visu}
    % \end{umlstate}




    \end{tikzpicture}
\end{document}

